\chapter{CONCLUSION AND FUTURE WORK}

% ─────────────────────────────────────────────────────────────────────────────
\section{Conclusion}

The aims of this thesis were to reduce the complexity of the H5P interactive video authoring workflow by designing a user-centred interface, and to make AI-assisted question generation viable for use in the production of interactive videos at a Vietnamese university. Both aims were achieved in a single web platform designed from scratch and tested against the dominant WordPress H5P workflow.

The core technical result is an architecture that integrates a React~18 + TypeScript SPA~\cite{react2023,typescript2023}, a stateless Express.js REST API~\cite{expressjs2023}, a SQLite data tier via Sequelize ORM~\cite{sqlite2023,sequelize2023}, and the Lumieducation headless H5P server library~\cite{lumieducation2023}. An AI pipeline extracts timestamped transcripts from video sources (YouTube captions, uploaded WebVTT files, or local Whisper transcription~\cite{radford2023}), passes them to a large language model (\texttt{llama-3.3-70b-versatile} via Groq, with a local Ollama fallback), validates the structured output with Zod schemas, and streams results to the browser via Server-Sent Events. Each suggestion carries a Bloom's taxonomy classification so instructors can gauge cognitive depth at a glance~\cite{anderson2001}.

The platform was evaluated against the WordPress H5P workflow in a comparative usability study with ten participants, each manually creating five randomly chosen interactions on the same video. Task completion rose from 30\% to 95\%; median task time fell from 18 minutes to 5 minutes; and the mean SUS score moved from 38.0 to 82.5. These improvements hold across every metric, which makes it difficult to attribute them to any single design choice. Putting all authoring tools on one screen, removing the need to switch between pages, and adding AI suggestion review together removed the barriers that made the WordPress workflow so slow.
\clearpage
\textbf{Revisiting the Research Questions.}

Full answers to each research question are developed in Chapter~\ref{ch:discussion}, Section~\ref{sec:rq_answers}. The key conclusions are summarised here.

RQ1 identified four principal usability barriers in the WordPress H5P workflow: multi-screen navigation fragmentation, CMS cognitive overhead irrelevant to the teaching task, absence of real-time preview, and no AI question scaffolding. All four were addressed in the new platform. RQ2 showed that co-locating all authoring components on a single viewport --- eliminating the split-attention effect~\cite{chandler1991} --- reduced the interaction count from 47 to 8 steps and drove most of the 13-minute task-time reduction. RQ3 established that a multi-provider LLM pipeline with Zod-based output validation and SSE streaming balances generation quality, perceived latency, and structural reliability. RQ4 found that task completion rose from 30\% to 95\%, median time-on-task fell from 18 minutes to 5 minutes, and the mean SUS score moved from 38.0 (``Poor'') to 82.5 (``Excellent'') on Brooke's scale~\cite{brooke1996}. These outcomes are consistent with Technology Acceptance Model predictions~\cite{davis1989} that reductions in perceived effort raise adoption likelihood, though confirming this causal link requires a longitudinal study.

\textbf{Summary of Contributions.}

Five contributions come out of this work:

\begin{enumerate}
  \item \textbf{A production-ready authoring platform:} An open-source, CMS-free H5P authoring environment designed for institutional deployment.
  \item \textbf{A functional AI-to-H5P pipeline:} The first documented integration of a production LLM pipeline directly generating complex, schema-compliant H5P content types guided by Bloom's taxonomy.
  \item \textbf{Structured output prompt engineering:} Reusable prompt matrices and Zod validation strategies that address the problem of forcing LLMs to produce valid nested instructional JSON.
  \item \textbf{Empirical evidence for UCD in EdTech:} Quantitative confirmation that a purpose-built user-centred interface substantially outperforms a general-purpose CMS workflow for this specific task.
  \item \textbf{Alignment with national digital education goals:} A locally deployable, open-source tool that supports VNU-HCM's digital transformation strategy without commercial licensing costs~\cite{vietnam_strategy2021}.
\end{enumerate}

\textbf{Limitations of the Study.}

Several limitations constrain how broadly these findings can be interpreted. The ten-participant sample satisfies Nielsen's guidance for identifying usability problems~\cite{nielsen1994}, but lacks the statistical power for parametric hypothesis testing. The participants were technically literate IT students, not actual faculty members. University professors may behave differently, particularly in how they respond to AI suggestions and recover from errors.

The AI quality evaluation used English-language transcripts. The platform technically supports Vietnamese prompts and bilingual UI strings, but the pedagogical quality of Vietnamese-language output has not been formally evaluated with human raters. This is a significant gap given the primary target institution. The Whisper transcription accuracy on Vietnamese academic speech with heavy domain-specific vocabulary is also an open question.

Finally, this evaluation measured short-term usability. Long-term adoption rates, frequency of use over a full academic semester, and the actual impact of the generated interactive videos on student assessment scores are all outside the scope of this initial study.

% ─────────────────────────────────────────────────────────────────────────────
\section{Future work}

The modular architecture opens several natural directions for extension.

\textbf{Longitudinal Faculty Evaluation Study.} A semester-long study with actual university faculty across diverse disciplines at VNU-HCM would provide the most ecologically valid evidence for the platform's impact. Measuring organic adoption rates, content production volume, and student performance data over a full academic term would answer questions this initial study could not.

\textbf{AI Provider Upgrade for Greater Concurrency.} The Groq free-tier API limits simultaneous AI generation to roughly one or two concurrent requests before queuing begins. Upgrading to a paid Groq tier, or switching the primary provider to a higher-capacity API such as Anthropic Claude or OpenAI GPT-4o, would allow the platform to serve a full classroom of instructors generating content simultaneously without rate-limit delays. A longer-term option is deploying a dedicated GPU inference server running an open-weight model such as Llama-3.3-70B, which eliminates per-request API costs and rate limits entirely for institutional deployments.

\textbf{Additional H5P Question Types.} The current platform supports six interaction types. Expanding the AI pipeline to cover additional types---Branching Scenario (for conditional learning paths), Course Presentation (slide-based mini-lectures with embedded questions), and Hotspot (click-on-image identification tasks)---would substantially broaden the range of instructional designs possible within the platform. Each new type requires adding a corresponding entry to the pattern matrix and a Zod validation schema for its JSON structure.
\clearpage
\textbf{Vietnamese Language Validation.} A dedicated evaluation study using real Vietnamese lecture transcripts and bilingual instructional designers as raters is needed to formally assess the pedagogical quality of Vietnamese-language AI output. Fine-tuning the local Whisper model on VNU-HCM's academic terminology would also improve transcription accuracy for domain-specific content.

\textbf{Deep Learning Analytics and xAPI Integration.} The current platform exports standard \texttt{.h5p} packages but does not capture learner interaction data. Implementing xAPI tracking that reports events (video replays, question failures, time per interaction) to a Learning Record Store would unlock learning analytics, allowing instructors to see which specific concepts students are struggling with~\cite{siemens2011}.

\textbf{Collaborative Real-Time Authoring.} Adding real-time collaborative editing would enable team-based content production between senior professors and teaching assistants. The backend already uses Express, so integrating WebSockets via \texttt{Socket.io} to broadcast state changes across connected clients is architecturally straightforward.

\textbf{Local LLM Fine-Tuning.} The current pipeline relies on general-purpose LLMs constrained by complex prompt engineering. Fine-tuning a smaller open-weight model (7B or 8B Llama variant) on a curated dataset of validated H5P question-transcript pairs from the VNU-HCM MOOC platform would improve suggestion quality and reduce inference latency for the local Ollama fallback path.

\textbf{Asynchronous Media Processing.} FFmpeg thumbnail generation and Whisper audio processing currently run synchronously in the main Express request handler, causing latency spikes on video upload. Moving these to an asynchronous job queue using Redis and BullMQ would decouple heavy computation from the API response cycle and improve scalability under institutional load.

\textbf{WCAG 2.1 AA Accessibility Audit.} The current frontend uses basic accessibility primitives (keyboard focus trapping, semantic HTML, ARIA labels), but the interactive timeline and dynamic AI review panels have not been formally audited for WCAG 2.1 AA compliance. Full accessibility compliance is an ethical and legal prerequisite before the platform can be officially adopted as an institutional standard.
